% PAGES: 268-268

We assume that it is possible to take arbitrarily large long or short positions in any
of the assets, and therefore the weights $w_i$ are not required to be positive or to be
smaller than 1.

Let $R$ be the return of the portfolio over a fixed period of time, and let $\mu_R =
E[R]$ be the expected value\footnote{Throughout this book, we also refer to the
expected value of the return of a portfolio as the ``expected return'' of the
portfolio.} of $R$. Let $R_i$ be the return (over the same period of time) of asset
$i$, and let $\mu_i = E[R_i]$ be the expected value of $R_i$, for $i = 1:n$. Let $r_f$
be the risk--free return of the cash position in the portfolio over the same period of
time.

The return of the portfolio is equal to the weighted average of the returns of the
assets and of the cash position, i.e.,
\begin{equation}
R \;=\; \sum_{i=1}^n w_iR_i \;+\; w_{cash}r_f;
\end{equation}
for completeness, a proof of (9.2) is given at the end of this section.

By taking expected values in (9.2) and using (9.1), we find that
\begin{align}
\mu_R &= \sum_{i=1}^n w_i\mu_i \;+\; w_{cash}r_f \notag\\
&= \sum_{i=1}^n w_i\mu_i \;+\; \left(1-\sum_{i=1}^n w_i\right)r_f \notag\\
&= r_f \;+\; \sum_{i=1}^n w_i(\mu_i-r_f).
\end{align}

Let $w = (w_i)_{i=1:n}$ be the vector of the asset weights in the portfolio, and let
$\mu = (\mu_i)_{i=1:n}$ be the vector of the expected values of the returns of the $n$
assets. Denote by $\mathbf{1}$ the $n\times 1$ column vector whose entries are all
equal to 1, and let
\begin{equation}
\bar\mu \;=\; (\bar\mu_i)_{i=1:n} \;=\; (\mu_i-r_f)_{i=1:n} \;=\; \mu-r_f\mathbf{1}
\end{equation}
be the $n\times1$ column vector of the excess expected returns of the $n$ assets over
the risk--free rate. Then,
\begin{equation}
\sum_{i=1}^n w_i\bar\mu_i \;=\; \bar\mu^tw,
\end{equation}
see (1.5), and (9.3) can be written in vector form using (9.5) as follows:
\begin{align}
\mu_R &= r_f \;+\; \sum_{i=1}^n w_i(\mu_i-r_f) \;=\; r_f \;+\; \sum_{i=1}^n
w_i\bar\mu_i \notag\\
&= r_f \;+\; \bar\mu^tw.
\end{align}

Also, note that (9.1) can be written in vector form using (1.5) as
\begin{equation}
w_{cash} \;=\; 1 - \mathbf{1}^tw.
\end{equation}

Let $\sigma_R$ be the standard deviation of the return $R$ of the portfolio, and let
$\sigma_i$ be the standard deviation of $R_i$, for $i = 1:n$. Let $\rho_{j,k}$ be the
correlation between the
